\section{Lean integration: a conservative bridge before an invasive one}
\label{sec:integration}
\status{Integration specification. The cited core names were inspected at tag
\texttt{v4.34.0}; every adapter interface below is proposed and none has been
compiled.}

\subsection{Begin at a stable boundary}

The first implementation should invoke \grind{} on ordinary Lean goals and
supply newly proved local facts explicitly. This is less efficient than an ideal
native integration, but it lets the structural planner and the certificate
pipeline be tested without depending immediately on every internal
representation.

There are concrete lower-level entry points to investigate next. In the
inspected v4.34.0 source, \code{Lean.Meta.Grind.Params} carries extensions,
extra theorems, extra facts, simplification data, and configuration;
\code{assertExtra} introduces supplied proof facts through
\code{addNewRawFact}; \code{mkGoalCore} initialises a goal and its solver
states; and the result structure carries failure information and
counters~\cite{grind-main}. \code{Lean/Meta/Tactic/Grind/Types.lean} exposes
\code{GrindM}, \code{GoalM}, opaque solver-extension state, and a
\code{SavedState} containing both Meta and grind state, whose restore operation
restores both components~\cite{grind-types}. There is also a registration path
for solver extensions, whose methods cover internalisation, new-equality and
new-disequality notification, an action, consistency checks, and a
model-combination hook~\cite{grind-api}.

A second inspection of the same release found a more specific surface: a
registration function \code{registerSolverExtension} in the
\code{Lean.Meta.Grind} namespace, together with
\code{SolverExtension.setMethods}, whose callbacks are \code{internalize},
\code{newEq}, \code{newDiseq}, \code{mbtc}, \code{action}, \code{check} and
\code{checkInv}; registration happens at initialisation time and extension state
is per goal. The API also distinguishes \code{Solvers.mkActionCore} from the
wrapper that drains pending facts~\cite{grind-api}. A polynomial adapter would
mark relevant expressions during internalisation, mark a component dirty after
an equality, perform bounded certificate search on the expensive-check path, and
publish its accepted output as a proved fact through the ordinary fact
machinery. Note that the internal-invariant check in that interface is a check
of data-structure bookkeeping; it is not the mathematical soundness argument for
a certificate.

These are implementation hooks, not a promise of a stable public integration
contract, and two qualifications matter.

First, feeding \code{Params.extraFacts} into a \emph{fresh} run is not the same
as resuming an earlier saturation state. A true incremental adapter must
preserve the goal state and all relevant solver and context state, and must add
facts through the normal internalisation and propagation path. Otherwise an
equality can exist in one table while being invisible to another solver.

Second, restoring a Lean metavariable snapshot alone is not enough.
\forge{} needs its own rollback layer for the additional clause-database views,
demand queues, proof dependencies, and branch-local plugin state that it
introduces. \code{SavedState} is necessary, not sufficient.

A proposed versioned adapter surface:

\begin{lstlisting}
GrindAdapter.open           (snapshot, options)     -> session
GrindAdapter.seed           (session, checkedProofs)-> session
GrindAdapter.advance        (session, workQuota)    -> localResult
GrindAdapter.exportRelevant (session, demand)       -> checkedSummary
GrindAdapter.snapshot       (session)               -> resumableState
GrindAdapter.restore        (resumableState)        -> session
\end{lstlisting}

All version-specific internals live in one compatibility module per pinned Lean
version --- for example \code{Forge.Bridge.GrindV434}. Theory adapters and the
outer scheduler target \forge{}'s own records rather than importing dozens of
evolving internal modules. Unit-test seeded facts, case-split rollback, proof
extraction, metavariable changes, and equality propagation through every enabled
local theory. Until that adapter is reliable, restarting a local tactic with
explicit checked facts is preferable to a fast but incoherent shared state.

\subsection{The tactic is the interface; it need not be the executor}
\label{sec:worker-split}

\S\ref{sec:integration} so far has assumed one axis of choice: how deeply the
tactic reaches into \grind{}'s internals. There is a second axis, and it is
arguably more consequential --- \emph{where the search runs}.

The obvious arrangement is that a native tactic performs the search in the Lean
server process. That has real advantages: exact access to the local context,
open namespaces, options and reducibility settings; direct return of an
\code{Expr} with no serialisation; natural integration with the language
server, tactic traces and source positions; automatic compilation against the
project's own toolchain; and invocation at an actual proof hole rather than at
a separate prompt. Any usable front end wants all of that.

But running the search there carries costs that a long-running planner feels
much more than an ordinary tactic does:

\begin{itemize}[leftmargin=1.4em]
\item a long or allocating search competes with elaboration and editor
  responsiveness;
\item cooperative cancellation only works if the search polls often enough;
\item a crash, an FFI failure, or a memory blow-up takes down the Lean server
  process;
\item keeping expensive search sessions alive \emph{across commands} requires
  environment extensions or external state whose lifecycle has to cooperate
  with incremental recompilation.
\end{itemize}

The last one deserves attention, because a persistent obligation graph spanning
a whole proof is exactly the thing this design wants and exactly the thing
incremental recompilation is entitled to discard.

The resolution is to keep the native tactic as the \emph{interface} while
moving execution into a project-local worker:

\begin{itemize}[leftmargin=1.4em]
\item a \textbf{controller} --- the tactic front end, or a standalone
  session --- owns user intent, settings, history, protocol framing, deadlines,
  process supervision and presentation. It holds \emph{no authority to declare
  a candidate well-typed, a goal unprovable, or a solver observation valid};
\item a \textbf{worker}, compiled under the target project's Lake environment,
  owns the \code{Environment}, elaboration, exact identities, reification,
  search, lowering to \code{Expr}, and kernel checking;
\item an optional \textbf{solver process} receives only a canonical query and
  returns bounded observations.
\end{itemize}

The tactic then prepares the exact goal, hands a \emph{pure prepared problem}
to the worker, and kernel-checks the returned expression in the server. A worker
crash loses process-local state only; the controller rebuilds from an import
root and a replay log.

This also answers a question \S\ref{sec:trust} raised and did not settle: where
the sandboxing boundary should be. It is not primarily about soundness --- an
exact certificate rejects a bad proposal wherever it was produced --- but about
not letting a numerical search take the editor with it.

\subsection{Three levels of identity}
\label{sec:identity}

\S\ref{sec:architecture} required scope keys to carry expression, type,
universe levels, instance identity, context fingerprint and scope, and warned
that printed text is not a safe key. That is right but under-specified: it
treats identity as one thing with several fields. It is better understood as
three distinct kinds of identity with different lifetimes and different
authority.

\begin{enumerate}[leftmargin=1.4em]
\item \textbf{Lean identity} --- exact \code{Name}, local declaration identity,
  environment generation, source \code{Expr}. Valid only inside one worker, one
  toolchain, one session generation.
\item \textbf{Neutral semantic identity} --- stable variable, global,
  constructor and provider IDs in a checked inventory. Serialisable,
  fingerprintable, independent of elaborator allocation.
\item \textbf{Presentation identity} --- delaborated syntax, qualified
  spelling, candidate ordinal, rendering route. \emph{Never sufficient to
  recover an authority.}
\end{enumerate}

A mapping from Lean identity to neutral identity is built during preparation
and kept in a private origin table; lowering uses the inverse mapping only
\emph{after} the neutral candidate has passed its checker; presentation is
derived last. The rule that falls out is sharper than anything in
\S\ref{sec:trust}'s protocol discussion:

\begin{invariant}[The origin table never leaves the worker]
\label{inv:origin}
An external process receives a sealed neutral problem and returns a neutral
candidate. No raw \code{Expr}, environment pointer, free-variable identifier,
process-local handle or certificate carrier crosses that boundary. Origin
resolution and kernel checking remain with the worker.
\end{invariant}

Environment generations make rollback safe in a way cache invalidation does
not. Attach a monotonically fresh generation token to every committed
environment state and index every cache, prepared origin, candidate token and
inventory by it. Undo then restores an older environment \emph{object} while
issuing a \emph{new} generation token, so stale capabilities cannot be reused
by accident --- which is a better answer to the stale-state failure mode of
\S\ref{sec:trust} than ``invalidate the affected caches'', because it fails
closed rather than depending on the invalidation being complete.

Durable snapshots, correspondingly, retain project and toolchain identity, root
imports, schema fingerprints, replay history and settings --- and retain no raw
\code{Expr}, no \code{FVarId}, no environment pointer, no process handle, no
candidate token, and no certificate carrier.

\subsection{Do not search in \texorpdfstring{\code{Expr}}{Expr}}
\label{sec:search-ir}

A related temptation is to drop the typed view entirely and search directly in
Lean's expression and metavariable machinery, on the grounds that this buys
definitional equality and typeclass synthesis for free. For a \emph{tactic}
step that is right. For a bounded, accounted, certificate-producing
\emph{search} it costs several properties this design depends on:

\begin{itemize}[leftmargin=1.4em]
\item a small, auditable fragment with an explicit completeness ledger
  (\S\ref{sec:negative});
\item stable alpha-normal fingerprints independent of elaborator allocation
  identifiers;
\item deterministic serialisation, and differential testing against an
  independent implementation;
\item independent checking, rather than replaying the search's own mutations;
\item exact accounting of steps, choices, pruning and plans --- without which
  \code{exhausted} cannot be distinguished from \code{truncated};
\item the ability to keep an unsupported dependent region opaque without it
  being normalised away under an aggressive reducibility setting;
\item decoupling from changes in Lean's metavariable implementation.
\end{itemize}

The metavariable context is excellent for final elaboration and for local
tactic work, and awkward as an authoritative branch representation: it is
optimised for elaborator workflows and is hard to serialise, fingerprint,
compare, replay and branch under an exact resource policy. The recommendation
is a Lean-oriented \emph{neutral} IR that may reference \code{Name}s, binder
kinds and exact declaration handles at the worker boundary, while owning its
own recursive structure and substitution rules --- and a pure branch-state
model for search, with \code{Expr} entering only at lowering.

\begin{remark}[Deduplication can discard the authority]
One consequence is worth spelling out because it is a silent failure. If
candidates are deduplicated by a lossy key --- pretty-printed text,
alpha-equivalence, eta, reducible unfolding, proof irrelevance --- a container
keyed by that equality can discard the representative that carries the checked
evidence and keep a look-alike that does not. Finish ranking and display
grouping over explicit keys, keep a separate pointer to the exact origin, and
do not define equality or ordering on the opaque authority merely for
convenience. Exact normalised neutral-tree identity is the safe default; text
remains useful as the user's observed grouping key and nothing more.
\end{remark}

\subsection{A typed semantic boundary for each worker}

A reifier does more than turn Lean expressions into strings. For a polynomial
fragment it produces a variable table, an exact polynomial, and a proof that the
polynomial's denotation under that table is the source expression:
\[
\operatorname{reify}(e)
  =\bigl(\text{environment }\rho,\ \text{IR }t,\
    \pi: e=\operatorname{eval}_{\rho}(t)\bigr).
\]
For an inequality or a formula the contract is an equivalence or a suitable
one-way implication, and the direction must match how the result will be used.
Unsupported subexpressions may become opaque atoms when that abstraction is
sound for the requested procedure.

An algebraically correct certificate for the wrong denotation is not a proof of
the original statement. Equality classes must be homogeneous up to justified
type compatibility; heterogeneous equality is not an excuse to ignore type
indices, so keep \code{HEq} reasoning and transport explicit and use elaboration
to validate casts. A proof of $i=j$ does not make \code{Vector A i} and
\code{Vector A j} definitionally interchangeable, and they must not be identified
because their indices were merged in an untyped side representation. Typeclass
dictionaries and universe instantiations are part of the elaborated expression,
not metadata that can always be erased.

An initial external term language should remain first-order and homogeneous
rather than implementing an ad hoc untyped approximation of all of Lean.

\subsection{Three levels of proof reconstruction}

\begin{enumerate}[leftmargin=1.4em]
\item \textbf{Ordinary theorem applications and tactic steps.} Appropriate for
  synthesised witnesses, explicit induction motives, small nonnegative linear
  combinations, and recurrence identities. The search can be complicated while
  the replay stays short.
\item \textbf{A typed certificate to a reflected checker}, with a theorem
  \[
  \operatorname{check}(P,C)=\texttt{true}\ \Longrightarrow\ \operatorname{Denotes}(P).
  \]
  The problem $P$ comes from the source reifier. The external solver may choose
  only $C$; it may never replace $P$. In a conservative configuration, reduce
  the checker inside the trusted proof-checking path rather than accepting an
  unchecked foreign result.
\item \textbf{External proof import} with explicit supported rules, every
  unsupported step becoming an obligation.
\end{enumerate}

For larger algebraic results the second level becomes worthwhile once its
soundness theorem exists; below that threshold, reconstructing small certificates
with \code{ring}, \code{linarith}, \code{positivity}, \code{norm\_num} and cast
normalisation avoids introducing a second independent arithmetic proof language
prematurely~\cite{mathlib-ring,mathlib-linarith,mathlib-positivity,mathlib-linear-combination}.

The mixed bit-vector case should use the existing shared-state
\code{grind => bv\_decide} route where applicable rather than duplicating its
encoding to fit a uniform architecture~\cite{lean-release-434}.

\subsection{A proposed module layout}

The following paths are an implementation plan. No file claimed here exists in
this repository.

\begin{lstlisting}
Forge/Frontend.lean              -- syntax, configuration, diagnostics
Forge/Core/GoalSnapshot.lean     -- restorable state identity
Forge/Core/Obligation.lean       -- AND/OR graph and proof assembly
Forge/Core/CheckedFact.lean      -- scoped propositions, dependencies
Forge/Core/Scheduler.lean        -- deterministic quotas and budgets
Forge/Core/Replay.lean           -- minimization and axiom audit

Forge/Bridge/GrindV434.lean      -- ALL version-pinned internals
Forge/Reify/Polynomial.lean      -- typed views and denotation proofs

Forge/Structural/CallAnalysis.lean
Forge/Structural/Generalize.lean
Forge/Structural/LemmaFrontier.lean
Forge/Structural/Witness.lean

Forge/Arithmetic/Cone.lean       -- cone semantics and exact replay
Forge/Arithmetic/Quadratic.lean
Forge/Arithmetic/Bernstein.lean
Forge/Arithmetic/Sturm.lean      -- needs a proved root-count theorem
Forge/Arithmetic/Lattice.lean

Forge/Instantiation/ConstrainedMatch.lean
Forge/Instantiation/ModelGuided.lean

Forge/Backend/SOS.lean           -- adapter to the existing sos project
Forge/Backend/LP.lean            -- adapter to the existing lp project
Forge/Backend/SMT.lean
Forge/Oracle/Protocol.lean       -- bounded certificate data, never code
Forge/Tactic.lean
\end{lstlisting}

The core modules depend only on Lean and a small shared representation.
Mathlib-specific reifiers and tactics sit behind separate imports. External
processes consume a versioned certificate format and return only untrusted
proposals, and the package remains useful when those processes are absent.

That last property is worth stating as a design goal rather than an accident:
the separation between core, checking, and external search permits
\emph{replay-only builds}. A numerical solver can be absent on the machine that
checks a finished proof.

\begin{longtable}{@{}>{\raggedright\arraybackslash}p{0.24\textwidth}>{\raggedright\arraybackslash}p{0.34\textwidth}>{\raggedright\arraybackslash}p{0.36\textwidth}@{}}
\toprule
Module & Responsibility & Acceptance contract \\
\midrule
\endhead
\code{Forge.Core} &
Scoped expression, fact, and obligation graph &
No unproved facts; valid dependency scopes \\
\code{Forge.Bridge} &
Version-specific extension hooks &
Restorable state; exact expression and proof transport \\
\code{Forge.Demand} &
Retrieval, instantiation, tabling &
Every application elaborates; acyclic proof parents \\
\code{Forge.Reify} &
Typed polynomial representation &
Denotation equality for each reified term \\
\code{Forge.Arithmetic} &
Cone, box, lattice certificate checking &
Proved checker soundness, or explicit proof reconstruction \\
\code{Forge.Structural} &
Recursor selection and generalisation &
Goal refinement through an actual induction theorem \\
\code{Forge.Synth} &
Invariants and existential candidates &
Base/step identities or witness specification discharged \\
\code{Forge.Backend} &
SAT, SMT, SOS, LP, superposition adapters &
Supported translation and proof rules; no raw oracle acceptance \\
\code{Forge.Replay} &
Proof minimisation and audit &
Exact original goal, no remaining metavariables, approved axioms \\
\bottomrule
\end{longtable}

\begin{proposition}[Compositional soundness is not inherited]
Soundness of the assembled system requires soundness at every layer
simultaneously. An unsound term translator is not repaired by an exact
polynomial checker. A sound checker is not repaired by a replay path that
assumes its result through an oracle axiom. An acyclic certificate DAG is not
sufficient if a proof parent was imported from the wrong branch.
\end{proposition}

\subsection{A proposed user interface}

Proposed syntax only; no such tactic exists.

\begin{lstlisting}
forge
forge? [selectedLemma, recursiveDefinition]
forge (config := {
  mode             := .compatibility
  induction        := true
  inductionDepth   := 2
  synthesizeLemmas := true
  witnesses        := .affine
  witnessSize      := 6
  polynomialDegree := 4
  externalSearch   := false
  trust            := .kernelOnly
})
\end{lstlisting}

A successful explanatory invocation emits a \emph{smaller replay plan}: the
chosen induction scheme, the generalised variables, the witness expressions, the
library lemmas actually needed, and the concrete certificate data. It must not
require the original broad search again, and it must not depend on an external
service still being available. A proof that succeeds only while a particular
model service is reachable is an unsatisfactory end product when an explicit
certificate can be retained instead.

Parameters must describe real implementation limits. A polynomial-degree bound
and a certificate monomial bound control different sources of explosion and
should not be conflated. A trust option must not change silently because a
faster worker happens to be available.

A failure report should identify the first substantive blocked obligation rather
than advising a larger timeout, and should distinguish at least:

\begin{itemize}[leftmargin=1.4em]
\item no induction motive found;
\item candidate lemma refuted by an executable example;
\item nonlinear certificate search exhausted;
\item unsupported source operation;
\item external proof reconstruction incomplete;
\item resource limit reached.
\end{itemize}

A mathematical failure and a resource failure must not produce the same
diagnostic. Useful concrete examples of the intended register: ``induction step
requires a lemma about the changing accumulator; 128 candidate expressions
rejected by tests; 3 universal proofs remain open,'' and ``SOS candidate
rejected by exact coefficient check; no theorem added.''

\subsection{The deployment pattern}

The intended end state for a user's file is not a permanently expensive tactic
call. It is: expensive search once, then a stable, readable proof for ordinary
recompilation. Concretely, the accumulator example should end as a separately
proved general lemma plus a short specialisation, with the final theorem not
invoking the synthesiser at all:

\begin{lstlisting}
theorem revAcc_spec (xs acc : List a) :
    revAcc xs acc = app (rev xs) acc := by
  induction xs generalizing acc with
  | nil => rfl
  | cons x xs ih =>
    simp only [revAcc, rev, ih, app_assoc, app]

theorem revAcc_correct (xs : List a) :
    revAcc xs [] = rev xs := by
  rw [revAcc_spec, app_right_nil]
\end{lstlisting}

The important algorithmic decision there is the motive --- \code{generalizing
acc} --- not the final normalisation tactic. A syntactically valid \code{induction}
call with a fixed accumulator gives an induction hypothesis too weak for the
recursive call, which is exactly the obstruction of \S\ref{sec:structural}.

\subsection{Implementation sequence and acceptance gates}

The nine contributing efforts proposed between four and six milestones each,
with substantial agreement on the ordering. The merged sequence is:

\paragraph{Gate 1 --- trustworthy orchestration.}
Frontend, deterministic compatibility mode, full state rollback, worker
contracts, replay records, and the axiom-policy audit. Exit criteria: the
restricted mode reproduces a pinned baseline suite without introducing new
axioms or losing source assumptions; a deliberately malformed proof or admitted
theorem is rejected; a failed branch provably cannot leak a local equality into
its sibling. Only existing proof-producing leaf tactics need run at this stage.
Status: a first milestone exists. The 	exttt{forge} tactic in
	exttt{Forge/Frontend.lean} tries a fixed sequence of workers, restores the whole
tactic state after each failure, and accepts a proof only after checking the term
for 	exttt{sorry} and for axioms outside the policy. Its tests pin a small
restricted baseline, reject an admitted proof, and show that a failed worker's
metavariable assignment does not survive; the last is tested, not proved, and
no planner exists.

\paragraph{Gate 2 --- exact algebraic replay.}
Typed reification and cone replay in Lean; the quadratic and finite-cone oracles
called through a bounded data protocol; the affine/Farkas and recurrence
certificate reconstructions. Exit criteria: successful Lean checking of every
claimed end-to-end example, mutation rejection \emph{at the data boundary}, and
recorded axiom dependencies. This stage already adds useful nonlinear facts to a
stock \grind{} leaf; it need not wait for an induction planner or SDP search.

\paragraph{Gate 3 --- structural planning and witnesses.}
Explicit theorem-application obligations, bounded induction with
dependency-directed generalisation, and typed witness grammars. Exit criteria:
unseen variations of accumulator correctness are solved \emph{with the
correctness theorem absent from the environment}; false-statement and
circular-lemma rejection tests pass; an induction success includes the selected
principle and reconstructed motive in its replay trace; a witness success
includes the witness term and every bound, integrality, and typeclass obligation
it needs.

\paragraph{Gate 4 --- domain-aware positivity.}
Bernstein coefficient and coverage rules formalised; compact-box certificates;
and, before any univariate certificate is accepted in Lean, a proved or
integrated Sturm/root-count result. That verification task is substantial and
must not be hidden behind an unchecked foreign call. An explicit multi-leaf
replay can validate the search-to-source path first, but production acceptance
should prefer a reusable checked certificate theorem to a large amount of
duplicated tactic text.

\paragraph{Gate 5 --- efficient cooperation.}
Incremental fact subscriptions, proof-DAG sharing, bound-derived product
envelopes, target-driven certificate candidates, and the version-pinned adapter
--- introduced only where measurement shows that repeated leaf setup dominates.
Keep a terminal-tactic fallback so that a Lean upgrade cannot disable the entire
system.

\paragraph{Gate 6 --- optional broader search.}
Premise search, SDP proposals with exact reconstruction, the CDCL(T) lane, and
optionally a learned ranking policy. Exit criterion for external backends: zero
unproved hints in any counted success. Learned components may choose which
candidate to try; they never decide which statements are accepted.
Native-trust workers remain opt-in, and every optional dependency carries its own
toolchain and feature-compatibility check.

\paragraph{What not to implement first.}
A complete new SMT solver, a rewrite of \grind{}, a general SDP backend, a
dependent-type-erasing universal intermediate language, or a learned scheduler.
The most valuable first release is not ``every solver behind one keyword''.

\subsection{Where the risk actually is}

The highest-risk engineering tasks are not mathematical. They are preserving
dependent contexts through generalisation; maintaining valid state across local
solver restarts and snapshots; producing compact proof reconstruction; and
keeping candidate lemmas out of the proved state. The numerical search is
comparatively easy to sandbox, precisely because an exact certificate can reject
a bad proposal.

A restricted, correct implementation that exports a small set of ordinary Lean
proofs is preferable to a wide solver bridge that silently leaves obligations or
enlarges the trusted base.
